\documentclass[12pt]{article}
\usepackage{fontspec}
\usepackage{lipsum}
\usepackage{polyglossia}
\setmainlanguage{french}
\usepackage{xcolor}
\usepackage{xspace}

\begin{document}
	
% Consignes (le texte est fourni; ajoutez ce qu'il faut dans le préambule) :
% 1. Ajouter un titre centré, en gras, en taille \Large.
% 2. Dans le premier paragraphe : un mot en emphase, un mot en petites capitales, un mot en rouge (\textcolor).
%    Que se passe-t-il si vous oubliez xcolor ? Lisez le .log.
% 3. Passer le deuxième paragraphe en \small (commande à bascule).
% 4. Le troisième paragraphe sans indentation(\noindent). 1ligner le titre et l'auteur à droite
% 5. Mettre un mot dans une \fcolorbox, puis un paragraphe entier dans l'environnement framed.
% Packages utiles : xcolor (couleurs, \fcolorbox), framed (environnement framed).
% 6. Pour aller plus loin : réutiliser une commande créée dans l'exo 2_5.



Le roman commence dans une gare de chemin de fer, une locomotive tonne, un postillon de piston couvre l’ouverture du chapitre, un nuage de fumée cache une partie du premier alinéa. Dans l’odeur de gare passe une bouffée d’odeur de buf‑ fet de gare. il y a quelqu’un qui regarde à travers les vitres embuées, il ouvre la porte vitrée du bar, tout est brumeux, même à l’intérieur, comme vu à travers les yeux d’un myope, ou à travers des yeux irrités par des escarbilles. ce sont les pages du livre qui sont embuées comme les vitres d’un vieux train, c’est sur les p  hrases que se pose le nuage de fumée. Il pleut ce soir‑là, l’homme entre dans le bar ; il déboutonne son pardessus humide ; un nuage de vapeur l’enveloppe ; un coup de sifflet s’en va au long des quais luisants de pluie à perte de vue. 

Un coup de sifflet qu’on dirait de locomotive et un jet de vapeur sortent de la machine à café que le vieux barman met sous pression comme s’il lançait un signal, ou c’est du moins ce qui ressort de la succession des phrases du second alinéa, où les joueurs attablés replient l’éventail de leurs cartes contre leur poitrine et se retournent vers le nouveau venu d’une triple torsion du cou, des épaules et des chaises, pendant que les clients au comptoir soulèvent les tasses et soufflent sur la surface du café les lèvres et les yeux entrouverts, ou sirotent la mousse des chopes de bière avec une attention excessive pour qu’elles ne débordent pas. le chat fait le dos rond, la caissière referme la caisse enregistreuse qui fait drin. l’ensemble de ces signes convergent vers une information : il s’agit d’une petite gare de province, où le nouvel arrivant se fait remarquer tout de suite.

Les gares se ressemblent toutes ; que les lampes ne parviennent pas à éclairer plus loin que leur halo importe peu, car c’est un endroit que tu connais par cœur, avec l’odeur de train qui reste même après que tous les trains sont partis, cette odeur spéciale des gares après que le dernier train est parti. tu as l’impression que les lampes de la gare et les phrases que tu es en train de lire ont pour tâche de dissoudre les choses qui affleurent d’un voile d’obscurité et de brouillard, plutôt que de les indiquer. Quant à moi, c’est la première fois de ma vie que je débarque dans cette gare, et il me semble déjà que j’y ai passé une vie entière, à entrer et à sortir de ce bar, à passer de l’odeur sous la véranda à l’odeur de sciure des cabinets, tout cela mélangé en une seule senteur qui est celle de l’attente, l’odeur des cabines téléphoniques quand il n’y a plus qu’à récupérer les jetons parce que le numéro appelé ne donne pas signe de vie.

Italo Calvino, Si par une nuit d'hiver un voyageur, chapitre 1, extrait

\end{document}